\chapter{The Christmas Tree Program}

The Christmas Tree is the principal programming layout. It is not a historical
spread and it is not merely a fixed tableau. It is a one-dimensional ordered
card stack unfolded into a bounded temporal structure.

\section{The geometry}

The six growing rows contain $1+2+3+4+5+6=21$ cards. A crown and a base make
22 cards, allowing the entire Major Arcana to occupy the structure exactly once
when the Major-only program is used.

The crown is the catalytic source that starts the clock. The six expanding rows
are the six further lifecycle steps. Crown plus rows therefore form a
seven-step deterministic lifecycle. The base is not an eighth time-step: it is
the terminal $T\_STOP$ boundary where the whole settles.

The crown and base are positionally entangled endpoints. This relation belongs
to the geometry, not to the particular cards placed there. The base contains
the tree as compressed possibility and terminal integration; the crown makes
that possibility present in phase-space.

\section{Programming with the tree}

The arrangement may be altered, prayed and meditated upon, then stored as a
reversible LIFO configuration. The Major positions supply the program grammar.
In a full-deck elaboration, Marseille Minor cards define local worldly details,
material conditions, relationships, resources, or other parameters inside the
Major architecture.

The same geometry can be used to probe a question, program a desired path, or
inspect a card from another reading. The position remains stable even when the
card occupying it changes.
